\chapter{Entropy Source: from Transistor to Min-Entropy}
\label{ch:ch3_entropy}

% [DRAFT v1 -- theory and simulation only; hardware figures pending]
% Source material: docs/modelo_ero_resultados.md, docs/arquitectura.md,
% analysis/ero_model.py, analysis/jitter_estimator.py, analysis/spice_ring.py

This chapter builds the chain that the rest of the thesis rests on: from
the physical jitter of a ring oscillator to a defensible figure of
min-entropy per output bit, and from that figure to the design parameters
of the sampler. Every step is validated before it is used. The
statistical estimators are validated against synthetic data of known
jitter; the analytic bounds are validated against an exact numerical
computation; and the transistor-level bench is validated by measuring its
own resolution before anything is measured with it. All figures in this
chapter are \simu{} unless stated otherwise; the hardware campaign is
reported in Chapter~\ref{ch:ch7_results}.

\section{The elementary ring oscillator}
\label{sec:ero}

The entropy source is an \ac{ERO}: a free-running ring oscillator RO1
whose output is sampled by a flip-flop clocked by a second ring
oscillator RO2 divided by an integer $K_D$~\cite{baudet_2011_security,
fischer_2014_embedded}. Between two samples the phase of RO1 accumulates
jitter for $K_D$ periods of RO2; when the accumulated jitter is
comparable to the period of RO1, the sampled bit becomes unpredictable.
The \ac{ERO} was chosen over higher-throughput alternatives (multi-ring
XOR, coherent sampling, transition-effect rings) for one reason: it is
the architecture for which a stochastic model with a provable entropy
bound exists and for which an embedded jitter measurement method has
been published and validated~\cite{petura_2016_survey}. Throughput is
not a requirement here --- the engine consumes entropy only to seed a
deterministic generator --- whereas a demonstrable bound is.

A design decision that follows from the platform is that \emph{the board
oscillator is never used as a reference}. The Basys~3 clock source is a
MEMS oscillator with a fractional \ac{PLL} whose datasheet specifies a
cycle-to-cycle jitter of up to 95~ps \est{}~\cite{digilent_2014_basys3rm},
one to two orders of magnitude above the per-period jitter of a ring
oscillator in a 28~nm process (2--4~ps \est{}~\cite{petura_2016_survey}).
Sampling RO1 with that clock would measure the phase noise of the PLL,
a global and potentially inducible component, rather than the thermal
noise of the ring. RO2 therefore samples RO1, and the 100~MHz system
clock only drives the synchronous logic.

\section{Stochastic model and the quality factor}
\label{sec:model}

Following Baudet \emph{et al.}~\cite{baudet_2011_security}, the phase
$\varphi(t)$ of RO1, measured in periods, is modeled as a Wiener process
with drift $\mu$ and volatility $\sigma^2$. Between two samples separated
by $\Delta t$ the phase advances $\nu = \mu\,\Delta t$ periods on average
and accumulates a variance
\begin{equation}
  Q = \sigma^2\,\Delta t ,
  \label{eq:Q}
\end{equation}
the \emph{quality factor}. The output bit is $s = 1$ when the fractional
phase lies in $[1/2, 1)$. Conditioned on the previous phase $\varphi_0$,
the bit is Bernoulli with
\begin{equation}
  P(s=1 \mid \varphi_0) = \sum_{k\in\mathbb{Z}}
  \left[\Phi\!\left(\tfrac{k+1-\varphi_0-\nu}{\sqrt{Q}}\right) -
        \Phi\!\left(\tfrac{k+\tfrac12-\varphi_0-\nu}{\sqrt{Q}}\right)\right],
  \label{eq:puno}
\end{equation}
where $\Phi$ is the standard normal distribution function; this is the
wrapped Gaussian evaluated exactly. Its first harmonic gives the maximum
bias $\varepsilon_{\max} \approx (2/\pi)\,e^{-2\pi^2 Q}$, from which
Baudet's Corollary~1 yields the well-known expression
\begin{equation}
  H(s \mid \varphi_0) = 1 - \frac{4}{\pi^2 \ln 2}\, e^{-4\pi^2 Q}
  + O\!\left(e^{-6\pi^2 Q}\right).
  \label{eq:baudet14}
\end{equation}

\section{Exact conditional entropy vs.\ the asymptotic bound}
\label{sec:exact}

Equation~\eqref{eq:baudet14} is routinely cited as a \emph{lower bound}
on entropy per bit. Read carefully, the inequality in the corollary is
only that the entropy rate of the sequence is at least the entropy of
one bit conditioned on the exact previous phase; the right-hand side is
an asymptotic expansion \emph{of} that conditional entropy, not a bound
on it. Expanding the binary entropy in the bias,
$h_2(\tfrac12+\varepsilon) = 1 - \tfrac{1}{\ln 2}\sum_{n\ge1}
\tfrac{(2\varepsilon)^{2n}}{2n(2n-1)}$, and averaging over a uniform
$\varphi_0$ with $\varepsilon = A\sin(2\pi\varphi)$,
$A = (2/\pi)e^{-2\pi^2 Q}$, gives
\begin{equation}
  H(s\mid\varphi_0) = 1 - \frac{4}{\pi^2\ln 2}\,e^{-4\pi^2 Q}
                        - \frac{8}{\pi^4\ln 2}\,e^{-8\pi^2 Q} - \cdots
  \label{eq:orden2}
\end{equation}
All omitted terms are negative, so \emph{any} truncation overestimates
the entropy; the next term decays as $e^{-8\pi^2 Q}$, not $e^{-6\pi^2 Q}$.
Table~\ref{tab:exact} compares the exact value computed from
\eqref{eq:puno} on a grid of 8192 phases with the first- and second-order
expansions. Adding the second-order term reduces the error of
\eqref{eq:baudet14} by a factor of 95 at $Q=0.10$ and 4970 at $Q=0.20$,
which confirms the derivation.

\begin{table}[htbp]
\centering
\caption{Conditional Shannon entropy per bit: exact wrapped-Gaussian
computation vs.\ truncated expansions \simu{}.}
\label{tab:exact}
\begin{tabular}{@{}lcccc@{}}
\toprule
$Q$ & Exact & Order 2 & Order 1 (Eq.~\ref{eq:baudet14}) & Error of order 1 \\
\midrule
0.02 & 0.7013194 & 0.7100952 & 0.7345213 & $+3.3\times10^{-2}$ \\
0.05 & 0.9163001 & 0.9164920 & 0.9187783 & $+2.5\times10^{-3}$ \\
0.10 & 0.9886728 & 0.9886733 & 0.9887174 & $+4.5\times10^{-5}$ \\
0.15 & 0.9984319 & 0.9984319 & 0.9984327 & $+8.5\times10^{-7}$ \\
0.20 & 0.9997823 & 0.9997823 & 0.9997823 & $+1.6\times10^{-8}$ \\
\bottomrule
\end{tabular}
\end{table}

In the operating region ($Q\approx0.2$) the discrepancy is $10^{-8}$ and
irrelevant; for an under-dimensioned generator ($Q\approx0.02$) the
truncated formula credits 3.3\,\% of entropy that does not exist. This
thesis therefore uses the exact computation as the conservative value and
cites \eqref{eq:baudet14} as what it is: an excellent asymptotic
approximation for $Q \gtrsim 0.1$. The same numerical machinery shows
that the min-entropy obtained from the first-harmonic bias,
$H_\infty = -\log_2(\tfrac12 + \varepsilon_{\max})$, agrees with the exact
worst-case value to better than $10^{-4}$~bit over $Q\in[0.05, 0.30]$.

\subsection{Design target}

AIS~20/31 version~3.0~\cite{peter_2024_ais31} admits two alternative
criteria for class PTG.2: Shannon entropy $\ge 0.9998$ or min-entropy
$\ge 0.98$ per internal bit. Solving both by bisection on the exact
entropies gives $Q \ge 0.2022$ and $Q \ge 0.2286$ respectively. The
min-entropy criterion is the more demanding one and is also the quantity
NIST SP~800-90B works with, so it is adopted here: \textbf{the design
target is $Q \ge 0.23$, with a set point of $Q = 0.30$} ($H_\infty =
0.9951$), the margin covering the reduction of thermal jitter at low
temperature, which is the standard adversarial test of Fischer and
Lubicz~\cite{fischer_2014_embedded}.

\section{Embedded jitter estimator and its two corrections}
\label{sec:estimator}

To choose $K_D$ one must know $\sigma^2$, and the method of Fischer and
Lubicz~\cite{fischer_2014_embedded} measures it from the bits themselves:
with $K_D = 1$, the fraction $c$ of pairs $(b_j, b_{j+M})$ that differ
within a block of $N$ pairs estimates the phase accumulated over $M$
samples, and its variance across $K$ blocks grows as $4 M Q_{\mathrm{raw}}$.
Implemented literally, with a straight-line fit, the estimator
\emph{underestimates} $Q_{\mathrm{raw}}$ by 13\,\% to 38\,\% depending on
the jitter level. Two causes were identified on synthetic data with known
$Q$ and corrected:

\begin{enumerate}
\item \textbf{Vertex effect.} $c$ is not linear in the accumulated phase;
it is a \emph{triangular} wave of period one and slope $\pm2$. When the
deterministic phase $M\nu \bmod 1$ falls near a vertex the distribution
straddles it and the variance collapses --- to 0.38 of its nominal value
in one measured case. The vertex is detected from the measured mean
itself ($\bar c \approx 0$ or $1$ is a vertex, $\bar c \approx \tfrac12$
the middle of the ramp), so distances $M$ with $\bar c$ outside
$[0.15, 0.85]$, or for which $3\sqrt{MQ}$ reaches the vertex, are
discarded iteratively.
\item \textbf{Non-constant estimation noise.} Estimating a fraction from
$N$ pairs adds a variance that scales with $\bar c(1-\bar c)$, which
changes with $M$. Fitting it as a constant intercept biases the slope; it
is fitted instead as a second regressor,
$\mathrm{Var}(c) = Q\,(4M) + \kappa\,\bar c(1-\bar c)/N$.
\end{enumerate}

After both corrections the estimator recovers the injected $Q_{\mathrm{raw}}$
within $\pm2\,\%$ over a range of $150\times$ (Table~\ref{tab:estimator}),
with a bias of $-0.6\,\%$ and a dispersion of 4.0\,\% over eight
independent seeds for $K=20\,000$ blocks. The fitted accumulation
exponent is $1.009$ for a simulator without flicker noise; measured on
hardware, the same exponent separates the thermal component ($b=1$, which
counts as entropy) from flicker and global jitter ($b=2$, which does not).

\begin{table}[htbp]
\centering
\caption{Recovery of a known $Q_{\mathrm{raw}}$ by the corrected embedded
estimator ($K=20\,000$ blocks of $N=100$ pairs) \simu{}.}
\label{tab:estimator}
\begin{tabular}{@{}cccc@{}}
\toprule
$Q_{\mathrm{raw}}$ injected & Estimated & Error & $M$ values used \\
\midrule
$2.0\times10^{-7}$ & $1.979\times10^{-7}$ & $-1.07\,\%$ & 19 \\
$5.0\times10^{-7}$ & $5.060\times10^{-7}$ & $+1.21\,\%$ & 18 \\
$2.0\times10^{-6}$ & $2.037\times10^{-6}$ & $+1.86\,\%$ & 16 \\
$8.0\times10^{-6}$ & $8.039\times10^{-6}$ & $+0.48\,\%$ & 9 \\
$3.0\times10^{-5}$ & $2.977\times10^{-5}$ & $-0.78\,\%$ & 5 \\
\bottomrule
\end{tabular}
\end{table}

A practical consequence for the hardware is that the blocks need not be
contiguous: each requires only $N+M$ consecutive bits, and independent
blocks give a slightly better estimate ($-0.30\,\%$ error with 20\,000
independent chunks). The capture buffer in the FPGA therefore holds a
single chunk of a few kilobits, filled at ring speed and read out slowly,
instead of megabits of contiguous samples.

\section{Transistor-level bench and its resolution}
\label{sec:spice}

The entropy source is the only block of the engine tied to the
technology, and in an integrated circuit it is designed at transistor
level. A SPICE bench was built around a five-stage ring with an enabling
NAND gate, with every technology-dependent parameter (device models,
widths, minimum length, supply) confined to a header so that moving from
the generic placeholder models used here to a foundry \ac{PDK} changes
one block of the netlist. Jitter is extracted from the threshold-crossing
instants of the output by the second difference
$D(M) = t_{k+2M} - 2t_{k+M} + t_k$, i.e.\ the Allan variance applied to
phase data, whose expectation $2M\sigma^2$ is free of the bias that a
simple difference over overlapping windows carries (measured at 5\,\%),
and which is the method used by Benea \emph{et al.} to characterize rings
both in Artix-7 and in 28~nm silicon~\cite{benea_2024_flicker}. Estimates
at different $M$ are combined with inverse-variance weights proportional
to the number of independent windows; without that weighting the
dispersion stayed at 15\,\% regardless of the simulated length.

The bench measures its own resolution first: with all noise sources
disabled, the residual jitter is the numerical noise of the simulator.
The first attempt gave 0.33~ps, of the same order as the jitter to be
measured; it was entirely the start-up transient of the oscillator, and
discarding the first twenty cycles brings the floor to \textbf{1.38~fs},
three orders of magnitude below the signal. Only then was noise injected.
Jitter scales with the square root of the injected spectral density
($\times2.01$ measured for $\times4$ in density, against $\times2.00$
expected) up to the point where the r.m.s.\ noise current reaches
about 80\,\% of the stage drive current; there the small-signal regime
breaks down, the mean period shifts, and the scaling law fails by a
factor of five. A validity criterion follows for every campaign,
including those on a real \ac{PDK}: keep injected noise below $\sim40\,\%$
of the stage current and verify that the mean period does not move when
noise is enabled. With the placeholder models the ring oscillates at
4.3~GHz, which is not physical; the numbers that enter the thesis come
from the foundry models, where transient noise analysis injects the
device noise itself.

\section{Design target and divider selection}
\label{sec:kd}

With $\sigma/T$ measured per period, $Q_{\mathrm{raw}} = (\sigma/T)^2$ and
the divider follows from $K_D = \lceil Q_{\mathrm{target}} /
Q_{\mathrm{raw}} \rceil$. For the range of jitter expected in a 28~nm
process the divider lies between $3\times10^4$ and $5\times10^5$ and the
throughput per ring between 1 and 20~kbit/s. The model reproduces the
published design point of Fischer and Lubicz ($K_D \approx 430\,000$ for
$\sigma = 5.01$~ps over $T = 7.81$~ns)~\cite{fischer_2014_embedded}. Such
rates are sufficient here: seeding the deterministic generator with
256~bits of full entropy requires about 512 raw bits at $H_\infty \approx
0.98$, i.e.\ between 0.03 and 0.5~s per seed, and the engine issues a seed
only at start-up and at reseed intervals. If more throughput were ever
needed, the clean route is several independent \acp{ERO}, each with its
own bound, rather than a multi-ring XOR whose independence cannot be
demonstrated~\cite{bochard_2010_true}.
